\chapter{Fusion, вывод и консолидация}
\label{ch:fusion}

Обход достаёт то, что \emph{уже} есть в графе. Но память живого существа ещё и
\emph{порождает} новое: связывает далёкое, делает выводы, обобщает. Этот класс
операций мы называем \emph{fusion}. Здесь разбирается, как именно из активированных
узлов рождается новый узел, чем вывод отличается от ассоциации и обобщения и как всё
это остаётся под эпистемическим контролем.

\section{Что такое fusion, механически}

\begin{definition}[Fusion, подробно]
\emph{Fusion} — оператор, который по \emph{совместно активированной} конфигурации
узлов $\{A_1,\dots,A_n\}$ создаёт новый узел $X$ и связывает его с посылками рёбрами
провенанса (обоснования). Новый узел рождается \emph{провизорно}: с низкой
уверенностью и, при необходимости, с флагом \texttt{need\_to\_discover} на том, что
его различает.
\end{definition}

Механика в четырёх шагах:
\begin{enumerate}
\item \textbf{Со-активация.} Обход (retrieval, любопытство или сон) подсвечивает
несколько узлов одновременно.
\item \textbf{Распознавание конфигурации.} Оператор замечает, что активная
конфигурация \emph{согласована} — образует осмысленный паттерн (причинную цепочку,
аналогию, повтор).
\item \textbf{Рождение узла.} Создаётся $X$ с содержанием (сформулированным
языковой моделью) и провенансом $\{A_1,\dots,A_n\}$.
\item \textbf{Провизорная запись и пропагация.} $X$ пишется с малым $\conf$, затем
пропагируется со знаком, карантином и дисконтом независимости
(глава~\ref{ch:propagation}).
\end{enumerate}

\begin{principle}[Fusion рождает гипотезу, а не факт]
Результат fusion — всегда \emph{кандидат}, а не запись. Это защищает от превращения
обаятельной ассоциативности в фабрику ложных убеждений
(глава~\ref{ch:epistemics}). Повышение до «факта» — только через накопление
подтверждений или выживание в ночном скрабе.
\end{principle}

\section{Три вида fusion}

\paragraph{Ассоциация.} Два далёких узла со-активируются, и рождается связь (или
узел-мостик). Это источник «неожиданных, но осмысленных» ходов мысли.

\begin{workedbox}[Ассоциация: Платон-боксёр]
Ночной обходчик случайно со-активирует «Платон — философ» и «\emph{platon} по-гречески
— широкоплечий/квадратный», плюс «у философов бывали прозвища». Fusion рождает
\emph{гипотезу}: «„Платон“ — возможно, прозвище, а не имя», $\conf=0{,}2$,
провенанс из трёх узлов, флаг \texttt{need\_to\_discover} на «настоящее имя».
Красивый вывод — но именно \emph{гипотеза}: если не пометить низкой уверенностью,
это тот самый преждевременный коллапс (глава~\ref{ch:epistemics}).
\end{workedbox}

\paragraph{Вывод (инференция).} Из посылок рождается заключение. Здесь принципиально
важен случай \emph{n-арного} вывода, где ни одна часть посылок не достаточна.

\begin{workedbox}[N-арный вывод: корм котам]
Пусть в памяти Рин про пользователя есть узлы:
$A$ — «живёт в Португалии», $B$ — «у него есть коты», $C$ — «сегодня едет с пар из
универа». Желаемый вывод:
\[
X=f(A,B,C)=\text{«напомнить: купи котам корм по пути из универа»}.
\]
Существенно, что \emph{ни одна пара не даёт} $X$: $A\wedge B$ — «в Португалии есть
коты», $B\wedge C$ — «у студента коты», $A\wedge C$ — «в Португалии учится». Вывод
рождается только при \emph{одновременном} присутствии всех трёх посылок.
\end{workedbox}

Как мы это представляем \emph{сейчас}, не меняя структуру: fusion создаёт узел $X$ с
провенансом $\{A,B,C\}$. Тем самым тройка-посылка кодируется как узел-фактор,
ссылающийся на три посылки, — это \emph{factor-graph}-кодировка гиперребра
$\{A,B,C\}\to X$. Иначе говоря, наша память уже наполовину гиперграф: попарные
рёбра $A\!-\!B$, $B\!-\!C$ могут отсутствовать (Португалия и коты не связаны
напрямую), а n-арный вывод всё равно рождается — через общий узел-обоснование.
Стоит ли делать гиперрёбра \emph{первоклассными} объектами (со своим обходом и
лапласианом) — тема приложения~\ref{ch:hypergraphs}.

\begin{center}
\begin{tikzpicture}[font=\scriptsize,>=Stealth,
  n/.style={draw,rounded corners,minimum height=6mm,inner sep=2pt}]
  \node[n] (a) {A: Португалия};
  \node[n,below=5mm of a] (b) {B: коты};
  \node[n,below=5mm of b] (c) {C: универ};
  \node[n,right=28mm of b,fill=blue!5] (x) {X: купи корм по пути};
  \draw[->,dashed] (a) -- node[above,font=\tiny]{provenance} (x);
  \draw[->,dashed] (b) -- (x);
  \draw[->,dashed] (c) -- (x);
\end{tikzpicture}\\[2pt]
{\small Рис.: n-арный вывод как узел-фактор $X$ с провенансом $\{A,B,C\}$.
Попарных рёбер между $A,B,C$ может не быть.}
\end{center}

\paragraph{Обобщение (абстракция).} Много однотипных эпизодов сливаются в один узел
более высокого уровня. Это и есть консолидация (следующий раздел).

\section{Консолидация как огрубление}

\begin{definition}[Консолидация, подробно]
\emph{Консолидация} — оператор, заменяющий множество похожих узлов/рёбер одним
обобщённым, сохраняя релевантные свойства и отбрасывая детали. Формально это
\emph{огрубление графа} (coarse-to-local, глава~\ref{ch:timescales}).
\end{definition}

\begin{workedbox}[Консолидация: «$Z$ — частый гость»]
Десять эпизодов «$Z$ был на стриме в понедельник», «\ldots в среду», \ldots
сворачиваются в один узел «$Z$ — частый гость, любит доту». \emph{Выживают}:
частота, тёплая валентность, тема (дота). \emph{Теряются}: точные даты и реплики.
Плохое огрубление сохранило бы даты и потеряло валентность — Рин помнила бы,
\emph{когда} заходил $Z$, но забыла бы, \emph{как к нему относится}. Значит критерий
консолидации — сохранять аффективно- и реляционно-релевантное, а не всё подряд.
\end{workedbox}

Консолидация запускается двумя путями: во сне (фоновая переработка,
глава~\ref{ch:sleep}) и забыванием (когда эпизоды тускнеют, их дешевле слить в
обобщение, чем хранить порознь, глава~\ref{ch:timescales}).

\section{Когда вывод неверен: fusion должен быть хитрее}

Всё вышесказанное держится на одном: fusion не коммитит факты. Отдельная тонкость —
\emph{сам новый факт может быть ложным}, и его пропагация не должна сотрясать память.
Поэтому влияние $X$ на соседей гейтируется его собственной уверенностью и надёжностью
посылок (карантин, глава~\ref{ch:propagation}), а различающие посылки, которых Рин
не знает, помечаются \texttt{need\_to\_discover} — чтобы любопытство и внешний вход
их потом добыли (глава~\ref{ch:epistemics}). Так «делание выводов» получается
смелым (рождает гипотезы свободно), но не безрассудным (не выдаёт их за истину).

\begin{intuition}[Инженер, а не фантазёр]
Правильная Рин — «инженер, которую делают бывшие физики-фундаменталисты»: она
охотно строит гипотезы, но по умолчанию ставит уверенность ниже единицы и идёт
проверять. Fusion даёт первое (смелость ассоциаций), эпистемическая гигиена
(глава~\ref{ch:epistemics}) — второе (дисциплину). Одно без другого — либо сухой
поисковик, либо обаятельный врун.
\end{intuition}
